\chapter{Lean as the Rendered Language}
\label{ch:lean}

\index{Lean}
The user-facing notation should render to a conservative subset of Lean. There is
no second opaque specification file: domain declarations are Lean definitions,
expected behaviours are Lean values, and requirements are Lean propositions or
checked queries.

\begin{lstlisting}
def checkout : Grammar CheckoutEvent :=
  request >>> (accepted <||> rejected)

def eventuallyAnswered : CTL Observation :=
  CTL.ag (CTL.implies requestPending
    (CTL.af (CTL.or requestAccepted requestRejected)))
\end{lstlisting}

Most semantic machinery should be imported from a fixed library. A system author
normally supplies only:

\begin{itemize}
  \item typed actors, identifiers, event kinds, and payloads;
  \item parameterised grammar definitions;
  \item projections from trace prefixes to observable propositions;
  \item named temporal requirements and proofs or decidable checks.
\end{itemize}

This boundary matters for automation. Generated Lean remains small enough to
review, while library proofs and algorithms are reused rather than regenerated.

\section{From specification to service}

Lean can host executable code, including a small web server. That makes the formal
artifact part of the running project rather than a document beside it. The same
typed definitions can drive validation endpoints, test-trace generation, protocol
schemas, visualisations, and runtime monitors. Extraction and integration choices
must preserve the distinction between a verified behavioural boundary and an
unverified infrastructure layer.
